\section{The Native RAA Relation Compiler}

We now restart Blaze at the model where the Blaze paper first lives: the native
multilinear IOP layer.  This is deliberately before Merkle trees, BaseFold,
systematic encodings, or concrete coordinate-opening implementations.  At this layer the prover
and verifier can name committed Boolean oracle handles \(R_h:\B^d\to\F\),
commit them to MLE oracle handles \(\widehat O_h:\F^d\to\F\), and ask for
field-point MLE evaluations of those hatted handles.

The goal of the native layer is not to prove proximity to a codeword.  Its goal
is to compile one algebraic statement into a small set of MLE evaluation
requests and Boolean-sum claims.  The later authentication layer will then make
those requests cryptographic.

\[
  \begin{aligned}
  \text{RAA relation}
  &\Longrightarrow
  (\EvalReqs,\SumClaims) \\
  &\Longrightarrow
  \text{batched sumcheck}
  \Longrightarrow
  \text{more }\EvalReqs .
  \end{aligned}
\]

\subsection{Native Relation}

Fix a field \(\F\).  Let \(K=2^k\) be the message length and \(N=2^n\) be the
RAA codeword length.  The base native relation is \(\RMLE[\PRAA]\).  The public
evaluation claim is
\[
  z\in\F^k,\qquad v\in\F,
\]
and the implicit input is an alleged codeword
\[
  y\in\F^{\B^n}.
\]
The prover witness is a message tensor
\[
  m\in\F^{\B^k}.
\]
The claim is that
\[
  y=\PRAA(m)
  \qquad\text{and}\qquad
  \MLE(m)(z)=v.
\]

The RAA map is public and has the stage form
\[
  \PRAA(m)
  =
  \Acc\circ M_{\pi_2}\circ\Acc\circ M_{\pi_1}\circ F_{\mathrm{rep}}(m),
\]
where \(F_{\mathrm{rep}}:\F^{\B^k}\to\F^{\B^n}\) repeats message entries,
\(M_{\pi_i}\) is the permutation operator for a public permutation
\(\pi_i:\B^n\to\B^n\), and \(\Acc\) is the prefix-sum map over
\(\B^n\).  We use the push-forward convention
\[
  (M_\pi x)[\pi(b)]=x[b]
  \qquad(b\in\B^n).
\]
Equivalently, \((M_\pi x)[a]=x[\pi^{-1}(a)]\).  This convention is chosen so
that the tagged-product permutation check below compares the pairs
\((x[b],\pi(b))\) with \((y[a],a)\).  The Blaze text states the permutation
test with the same tagged leaves but writes the relation in pullback form
\(y[b]=x[\pi(b)]\)~[Blaze24].  In this explainer we fix the convention that
makes the product argument literal; if one instead uses the pullback
convention, the same checker should be read with \(\pi^{-1}\) in place of
\(\pi\).

Thus an honest prover's stage tensors satisfy
\[
  \begin{aligned}
  u_0&=m,&
  u_1&=F_{\mathrm{rep}}(u_0),&
  u_2&=M_{\pi_1}u_1,&
  u_3&=\Acc(u_2),\\
  u_4&=M_{\pi_2}u_3,&
  u_5&=\Acc(u_4)=y.
  \end{aligned}
\]

The prover provides Boolean oracle handles for
\[
  u_0=m\in\F^{\B^k},
  \qquad
  u_2,u_3,u_4\in\F^{\B^n},
  \qquad
  f_1,g_1,f_2,g_2\in\F^{\B^{n+1}}.
\]
The native MLIOP functionality commits those Boolean handles into hatted MLE
handles before the verifier asks field-point evaluation questions.
The prover does not send \(u_1\) or \(u_5\) as separate native messages.
Queries to \(u_1\) are emulated from \(u_0\) using the repetition layout:
we write
\[
  \widehat O_{u_1}:=\RepView(\widehat O_{u_0}),
  \qquad
  \widehat O_{u_1}(r):=\widehat O_{u_0}(\mathsf{msg}(r)).
\]
This is a derived view, not a fresh call to
\(\mathcal F_{\mathrm{MLIOP}}.\mathsf{Commit}\).  The final stage \(u_5\) is
the input codeword handle supplied to the checker.

\subsection{Request And Claim Sets}

A native evaluation request is a tuple
\[
  (J,r,v)\in\EvalReq,
\]
where \(J\) is a handle that can later authenticate MLE values, \(r\in\F^d\)
has the handle's dimension, and \(v\in\F\).  In the native MLIOP section \(J\)
is a hatted MLE handle \(\widehat O:\F^d\to\F\).  In the compiled Blaze section
the same slot is a coordinate-opening handle \(C\).  The request means that the
later verifier check needs
\[
  J(r)=v.
\]
For example, \((\widehat O_{u_3},r,s)\) means that the value \(s\) must be the
MLE evaluation of the tensor \(u_3\) at \(r\).

A Boolean-sum claim is stored in \(\SumClaims\).  Its logical content is
\[
  (s,G)\in\SumClaims
  \quad\Longleftrightarrow\quad
  \sum_{X\in\B^d}G(X)=s.
\]
In protocol boxes below, the reducers do not return an opaque polynomial name.
They return a small record that says which public residual form is being used,
which MLE handle it touches, which projection point was sampled, and what the
Boolean-sum target is.  The batched sumcheck then expands those records
explicitly.

The native layer uses two kinds of outputs:
\[
  \EvalReqs,\qquad \SumClaims.
\]
\(\EvalReqs\) are MLE claims to authenticate later.  \(\SumClaims\) are
Boolean sums to feed into sumcheck.  Some apparent scalar equalities, such as
the permutation root equality, are encoded by reusing or deriving values inside
\(\EvalReqs\) instead of carrying a separate equality set.

\subsection{How Native Tensor Checks Compose}

The same pattern appears in every RAA stage.  Start with a large relation
between named tensors.  Rewrite that relation as a residual tensor
\[
  R\in\F^{\B^s}
\]
that should be zero at every Boolean coordinate:
\[
  R[a]=0
  \qquad(a\in\B^s).
\]
Each residual value is a public low-degree expression in MLE evaluations of the
named tensors.  Abstractly,
\[
  R[a]
  =
  \Phi_a\Bigl(
    \MLE(t_{j_1})(\psi_1(a)),
    \ldots,
    \MLE(t_{j_q})(\psi_q(a))
  \Bigr),
\]
where \(t_{j_1},\ldots,t_{j_q}\) are the participating tensors, the address
maps \(\psi_t:\B^s\to\F^{d_{j_t}}\) and the expression \(\Phi_a:\F^q\to\F\)
are public.  The subscript \(a\) only means that the coefficients of the
expression may depend on the residual coordinate.
To use sumcheck, the Boolean-coordinate formula must come with a low-degree
field extension: the maps \(\psi_t\) and the expression \(\Phi_a\) are extended
to field points in the way specified by the concrete check.  This is not an
extra commitment.  It is part of the algebraic description of the residual
identity being checked.

The verifier does not read all of \(R\).  It samples a field point
\[
  r\in\F^s
\]
and projects the zero claim to one scalar equation:
\[
  \sum_{a\in\B^s}\chi_a(r)R[a]=0.
\]
If \(R\) is not the zero tensor, this random projection fails with high
probability over \(r\).  When the projected sum is still too large to compute
directly, the prover and verifier run sumcheck on the projected identity.
Sumcheck does not authenticate any tensor value by itself; it only reduces the
identity to endpoint claims of the form
\[
  \MLE(H_j)(r_j)=v_j.
\]
Those endpoint claims become \(\EvalReqs\).  The native layer records them; the
later authentication layer checks them.

A small example is the pointwise product relation
\[
  C[a]=A[a]B[a]\qquad(a\in\B^s).
\]
The residual is
\[
  R[a]=C[a]-A[a]B[a].
\]
The projected zero claim is
\[
  \sum_{a\in\B^s}\chi_a(r)\bigl(C[a]-A[a]B[a]\bigr)=0.
\]
If the prover cheats at one coordinate \(a_0\), the bad contribution is weighted
by \(\chi_{a_0}(r)\).  The point is not that the verifier checks the coordinate
\(a_0\) directly.  The point is that a random linear projection turns a large
coordinate-wise relation into one low-degree scalar identity, and sumcheck then
turns that identity into a few MLE endpoint requests.

\subsection{Accumulator Reducer}

The accumulator relation has the form
\[
  y=\Acc(x),
  \qquad
  x,y\in\F^{\B^n}.
\]
Let
\[
  \widehat A:\F^n\times\F^n\to\F
\]
be the multilinear extension of the accumulator matrix, where the Boolean
matrix has entry \(A[b,a]=1\) exactly when \(\Int(b)\ge \Int(a)\) in the
natural binary order.  In ordinary coordinate language,
\[
  y[b]=\sum_{a\le b}x[a].
\]
This matrix is conceptually \(N\times N\), but the verifier never materializes
it.  For fixed \(r\in\F^n\), the weights
\((\widehat A(r,X))_{X\in\B^n}\) can be generated in one linear pass over the
Boolean domain, and a single field-point value
\(\widehat A(r,s)\) can be computed by the usual recursive comparison formula
for the predicate \(\Int(b)\ge\Int(a)\).  Thus the reducer uses the
accumulator matrix as a compact algebraic kernel, not as a dense table.

The verifier cannot check this equation for every \(b\).  Instead it samples a
field point \(r\in\F^n\) and asks for the same linear relation at that MLE
point.  For an honest accumulator,
\[
  \MLE(y)(r)
  =
  \sum_{X\in\B^n}\widehat A(r,X)\,x[X]
  \qquad
  \text{for every }r\in\F^n.
\]
The value \(\MLE(y)(r)\) is not trusted for free.  The prover sends it as a
scalar \(s\), the reducer records \((\widehat O_y,r,s)\) as an evaluation request,
and the Boolean-sum claim says that the accumulator input \(x\) must sum to the
same scalar \(s\).

\begin{protocolbox}{Subroutine: Native Accumulator Reduction}{fig:native-accumulate-reduce}
\scriptsize\raggedright
\noindent Public parameters: \(n\), \(\widehat A:\F^n\times\F^n\to\F\).

\smallskip
\noindent\(\underline{\AccReduce(\verifier:(\widehat O_x:\F^n\to\F,\ \widehat O_y:\F^n\to\F,\ r\in\F^n),\ \prover:(x,y\in\F^{\B^n}))\to(\EvalReqs,B_{\mathrm{acc}})}\)
\begin{enumerate}[leftmargin=0.24in,label=\textbf{A\arabic*.},itemsep=0.06em,topsep=0.06em,parsep=0pt]
  \item \(\prover\): \(s_y\gets \MLE(y)(r)\).
  \item \(\prover\to\verifier\): \(s_y\in\F\).
  \item \(\verifier\): \(\EvalReqs\gets\{(\widehat O_y,r,s_y)\}\).
  \item \(\verifier\): \(B_{\mathrm{acc}}\gets(\mathsf{acc},s_y,\widehat O_x,r)\).
  \item \(\verifier\): return \((\EvalReqs,B_{\mathrm{acc}})\).
\end{enumerate}
\end{protocolbox}

The record \(B_{\mathrm{acc}}=(\mathsf{acc},s_y,\widehat O_x,r)\) means
\[
  \sum_{X\in\B^n}\widehat A(r,X)\MLE(x)(X)=s_y,
\]
where \(\widehat O_x\) is the handle that will later authenticate endpoint
values of \(\MLE(x)\).  The scalar \(s_y\) is local to the reducer.  It is not a
separate output because it is already stored in the evaluation request and in
the Boolean-sum target.
In the RAA relation, this reducer is used twice:
\[
  (u_2,u_3)
  \qquad\text{and}\qquad
  (u_4,y).
\]

\subsection{Permutation Reducer}

The permutation relation has the form
\[
  y=M_\pi x,
  \qquad
  y[\pi(b)]=x[b]\quad\text{for every }b\in\B^n.
\]
The verifier samples \(\alpha,\beta\in\F\).  The prover then constructs two
product-tree tensors \(f,g\in\F^{\B^{n+1}}\).  The leaves are intended to be
\[
  f[\mathsf{leaf}(b)]
  =
  \alpha-\bigl(x[b]+\beta\,\Idx(\pi(b))\bigr),
\]
\[
  g[\mathsf{leaf}(b)]
  =
  \alpha-\bigl(y[b]+\beta\,\Idx(b)\bigr).
\]
If \(y[\pi(b)]=x[b]\) and \(\pi\) is a bijection, these two leaf multisets have
the same product.  The product trees make this global product comparison
locally checkable.

Here is the intuition.  The left product writes each input value together with
the address where it is supposed to land:
\[
  \prod_{b\in\B^n}
    \left(\alpha-\bigl(x[b]+\beta\,\Idx(\pi(b))\bigr)\right).
\]
The right product writes each output value together with its actual address:
\[
  \prod_{a\in\B^n}
    \left(\alpha-\bigl(y[a]+\beta\,\Idx(a)\bigr)\right).
\]
If \(y[\pi(b)]=x[b]\), then renaming the right product by \(a=\pi(b)\) gives
\[
  y[a]=y[\pi(b)]=x[b],
  \qquad
  \Idx(a)=\Idx(\pi(b)),
\]
so the two products are term-by-term equal after reindexing.  If the tagged
multisets differ, the two products agree only with small probability over the
random \(\alpha,\beta\).  The address tag \(\beta\,\Idx(\cdot)\) matters: it
prevents the prover from merely preserving the multiset of values while moving
them to the wrong addresses.

We use the source product-tree coordinate convention
\[
  \mathsf{root}=(1^n,0),\qquad
  \mathsf{leaf}(b)=(0,b),
\]
\[
  \mathsf{parent}(X)=(1,X),\qquad
  \mathsf{left}(X)=(X,0),\qquad
  \mathsf{right}(X)=(X,1).
\]
The same symbols are also the low-degree field maps used in the sumcheck
formulas:
\[
  \mathsf{leaf}(r)=(0,r),\quad
  \mathsf{parent}(r)=(1,r),\quad
  \mathsf{left}(r)=(r,0),\quad
  \mathsf{right}(r)=(r,1)
  \qquad(r\in\F^n).
\]
For \(T\in\F^{\B^{n+1}}\), define
\[
  \mathsf{TreeRes}(T,X)
  =
  \widehat T(\mathsf{parent}(X))
  -
  \widehat T(\mathsf{left}(X))\,\widehat T(\mathsf{right}(X)).
\]
On Boolean \(X\), the equation \(\mathsf{TreeRes}(T,X)=0\) says that the node
named by \(\mathsf{parent}(X)\) is the product of the two children named by
\(\mathsf{left}(X)\) and \(\mathsf{right}(X)\).  Therefore a valid product tree
has a root equal to the product of all its leaves.  The coordinate
\(1^{n+1}\) is the padding point above the root; the honest construction sets
it to \(0\), so the residual at \(X=1^n\) also vanishes.  In protocol boxes,
\(\mathsf{TreeRes}(H_T,X)\) means this same expression with tensor label
\(H_T\); the prover instantiates the polynomial during sumcheck, and the
verifier records the required endpoint evaluations.
For a non-Boolean sumcheck endpoint \(q\in\F^n\), this is a formal-extension
evaluation, not a coordinate opening of \(T\).  Concretely, evaluating
\(\mathsf{TreeRes}(T,q)\) requires the three MLE values
\[
  \widehat T(\mathsf{parent}(q)),\qquad
  \widehat T(\mathsf{left}(q)),\qquad
  \widehat T(\mathsf{right}(q)).
\]
The Blaze prose later says that this contributes two queries per product-tree
tensor.  With the residual written above, the scalar-tensor realization needs
three formal-extension values.  We therefore treat the two-query line as a typo
or as accounting for a virtual/packed object not specified there.  In this
paper all three evaluation claims are made explicit.

\begin{protocolbox}{Subroutine: Native Permutation Reduction}{fig:native-permute-reduce}
\scriptsize\raggedright
\noindent Public parameters: \(n\), \(\Idx,\widehat\pi:\F^n\to\F\), product-tree maps \(\mathsf{root},\mathsf{leaf},\mathsf{parent},\mathsf{left},\mathsf{right}\), \(\mathcal F_{\mathrm{MsgOracle}}\), and \(\mathcal F_{\mathrm{MLIOP}}\).

\smallskip
\noindent\(\underline{\PermuteCommit(\verifier:(\widehat O_x,\widehat O_y,\alpha,\beta),\ \prover:(x,y))\to(\verifier:(\widehat O_f,\widehat O_g),\prover:(\widehat O_f,\widehat O_g,f,g))}\)
\begin{enumerate}[leftmargin=0.24in,label=\textbf{P\arabic*.},itemsep=0.06em,topsep=0.06em,parsep=0pt]
  \item \(\prover\): for every \(b\in\B^n\), let
        \(f[0\Vert b]\gets\alpha-(x[b]+\beta\,\widehat\pi(b))\).
  \item \(\prover\): for every \(b\in\B^n\), let
        \(g[0\Vert b]\gets\alpha-(y[b]+\beta\,\Idx(b))\).
  \item \(\prover\): for \(s\in\{1,\ldots,n\}\) and
        \(a\in\B^{n-s}\), let
        \(f[1^s\Vert0\Vert a]\gets
        f[1^{s-1}\Vert0\Vert a\Vert0]\,
        f[1^{s-1}\Vert0\Vert a\Vert1]\).
  \item \(\prover\): for \(s\in\{1,\ldots,n\}\) and
        \(a\in\B^{n-s}\), let
        \(g[1^s\Vert0\Vert a]\gets
        g[1^{s-1}\Vert0\Vert a\Vert0]\,
        g[1^{s-1}\Vert0\Vert a\Vert1]\).
  \item \(\prover\): let \(f[1^{n+1}]\gets0\) and
        \(g[1^{n+1}]\gets0\).
  \item \(\prover,\verifier\): \(R_f\gets
        \mathcal F_{\mathrm{MsgOracle}}.\mathsf{Submit}(\prover:f,\verifier:\bot)\).
  \item \(\prover,\verifier\): \(R_g\gets
        \mathcal F_{\mathrm{MsgOracle}}.\mathsf{Submit}(\prover:g,\verifier:\bot)\).
  \item \(\prover,\verifier\): \(\widehat O_f\gets
        \mathcal F_{\mathrm{MLIOP}}.\mathsf{Commit}(\prover,\verifier:R_f)\).
  \item \(\prover,\verifier\): \(\widehat O_g\gets
        \mathcal F_{\mathrm{MLIOP}}.\mathsf{Commit}(\prover,\verifier:R_g)\).
  \item \(\prover,\verifier\): return \((\verifier:(\widehat O_f,\widehat O_g),\prover:(\widehat O_f,\widehat O_g,f,g))\).
\end{enumerate}

\smallskip
\noindent Types: \(\widehat O_x,\widehat O_y,\widehat\pi,\Idx:\F^n\to\F\);
\(\widehat O_f,\widehat O_g:\F^{n+1}\to\F\);
\(x,y\in\F^{\B^n}\); \(f,g\in\F^{\B^{n+1}}\);
\(\alpha,\beta\in\F\); \(r\in\F^n\).

\smallskip
\noindent\(\underline{\PermuteReduce(\verifier:(\widehat O_x,\widehat O_y,\widehat O_f,\widehat O_g,\widehat\pi,\alpha,\beta,r),\ \prover:(x,y,f,g))\to(\EvalReqs,B_f,B_g)}\)
\begin{enumerate}[leftmargin=0.24in,label=\textbf{R\arabic*.},itemsep=0.06em,topsep=0.06em,parsep=0pt]
  \item \(\prover\): \(v_{\mathrm{root}}\gets \MLE(f)(\mathsf{root})\).
  \item \(\prover\): \(v_x\gets \MLE(x)(r)\) and \(v_y\gets \MLE(y)(r)\).
  \item \(\prover\to\verifier\): \(v_{\mathrm{root}},v_x,v_y\in\F\).
  \item \(\verifier\): \(v_f\gets\alpha-(v_x+\beta\,\widehat\pi(r))\) and
        \(v_g\gets\alpha-(v_y+\beta\,\Idx(r))\).
  \item \(\verifier\): \(\EvalReqs\gets
        \{(\widehat O_f,\mathsf{root},v_{\mathrm{root}}),
        (\widehat O_g,\mathsf{root},v_{\mathrm{root}}),
        (\widehat O_x,r,v_x),(\widehat O_y,r,v_y),
        (\widehat O_f,\mathsf{leaf}(r),v_f),
        (\widehat O_g,\mathsf{leaf}(r),v_g)\}\).
  \item \(\verifier\): \(B_f\gets(\mathsf{tree},0,\widehat O_f,r)\).
  \item \(\verifier\): \(B_g\gets(\mathsf{tree},0,\widehat O_g,r)\).
  \item \(\verifier\): return \((\EvalReqs,B_f,B_g)\).
\end{enumerate}
\end{protocolbox}

The records \(B_f=(\mathsf{tree},0,\widehat O_f,r)\) and
\(B_g=(\mathsf{tree},0,\widehat O_g,r)\) mean
\[
  \sum_{X\in\B^n}\chi_r(X)\mathsf{TreeRes}(f,X)=0,
  \qquad
  \sum_{X\in\B^n}\chi_r(X)\mathsf{TreeRes}(g,X)=0.
\]
The handle stored in each record tells the later endpoint checks which
committed product-tree tensor the values must come from.

The RAA relation uses this reducer twice:
\[
  (u_1,u_2,\pi_1)
  \qquad\text{and}\qquad
  (u_3,u_4,\pi_2).
\]
Therefore it introduces four product-tree tensors
\[
  f_1,g_1,f_2,g_2\in\F^{\B^{n+1}}.
\]
Each tree has \(2^{n+1}=2N\) entries.

For each permutation edge, read the checks as one correlated argument:
the derived leaf values load the claimed products from \(x\) and \(y\), the
parent-child \(\SumClaims\) force \(f\) and \(g\) to be product trees, and the
shared root value compares the two products.  No single one of these checks
proves a permutation by itself.

The top-level Blaze box below performs the native RAA orchestration directly.
It creates the stage handles, calls the two reducers, gathers their returned
requests, and runs one batched sumcheck over the six Boolean-sum claims in a
fixed order.  The \(u_1\) handle is virtual, but not opaque: its evaluation
rule is written where it is created.

\subsection{The Batched Six-Sumcheck}

The two accumulator reducers contribute two Boolean-sum claims \(B_{u2}\) and
\(B_{u4}\).  The two permutation reducers contribute four product-tree
Boolean-sum claims \(B_{f1},B_{g1},B_{f2},B_{g2}\).  Blaze batches these six
named claims into
one sumcheck by sampling one scalar
\[
  \lambda\leftarrow\F
\]
and using the powers \(1,\lambda,\lambda^2,\ldots,\lambda^5\) as batching
weights.
This is only native batching of Boolean-sum claims.  It is not the later
compiler-level batching of arbitrary MLE evaluation requests.  If any one of
the six Boolean-sum identities is false, then the vector of six errors defines
a nonzero polynomial in \(\lambda\) of degree at most \(5\).  The batching step
can hide that error with probability at most \(5/|\F|\).  Conditioned on the
batched claim being false, ordinary sumcheck soundness applies.  This is the
same scalar-power batching pattern used in Blaze's Remark~2.3.

Only the Boolean-sum identities are batched here.  The permutation leaf
equations and root equality have already been encoded as evaluation requests by
deriving leaf values and reusing one root value.  The endpoint values produced
by this sumcheck are still not trusted; they are appended to \(\EvalReqs\) and
authenticated later by \(\MLEAuth\).
This is exactly the endpoint discipline from
Figure~\ref{fig:subroutine-sumcheck-terminal}: \(\SumCheck_n\) produces one
terminal scalar \(T_{\mathrm{sc}}\), and the surrounding RAA flow expands that scalar
into MLE requests against the existing handles \(\widehat O_{u_2}\),
\(\widehat O_{u_4}\), \(\widehat O_{f_1}\), \(\widehat O_{g_1}\),
\(\widehat O_{f_2}\), and \(\widehat O_{g_2}\).

\begin{protocolbox}{Subroutine: Native Six-Sumcheck}{fig:native-six-sumcheck}
\scriptsize\raggedright
\noindent Public parameters: \(n\), \(\widehat A\), \(\chi_r\), \(\mathsf{TreeRes}\), and \(\SumCheck_n\) from Figure~\ref{fig:subroutine-sumcheck-terminal} with individual degree bound \(2\).

\smallskip
\noindent\(\underline{\RAASixSumcheck(\verifier:(B_{u2},B_{u4},B_{f1},B_{g1},B_{f2},B_{g2}),\ \prover:(u_2,u_4\in\F^{\B^n}, f_1,g_1,f_2,g_2\in\F^{\B^{n+1}}))\to(a_{\mathrm{sc}}\in\{0,1\},\EvalReqs_{\mathrm{sc}})}\)

\smallskip
\begin{enumerate}[leftmargin=0.24in,label=\textbf{S\arabic*.},itemsep=0.06em,topsep=0.06em,parsep=0pt]
  \item \(\verifier\): let
        \(B_{u2}=(\mathsf{acc},s_{u2},\widehat O_{u_2},r_{u2})\) and
        \(B_{u4}=(\mathsf{acc},s_{u4},\widehat O_{u_4},r_{u4})\).
  \item \(\verifier\): let
        \(B_{f1}=(\mathsf{tree},0,\widehat O_{f_1},r_{f1})\),
        \(B_{g1}=(\mathsf{tree},0,\widehat O_{g_1},r_{g1})\),
        \(B_{f2}=(\mathsf{tree},0,\widehat O_{f_2},r_{f2})\), and
        \(B_{g2}=(\mathsf{tree},0,\widehat O_{g_2},r_{g2})\).
  \item \(\verifier\): sample \(\lambda\leftarrow\F\), then send
        \(\lambda\) to \(\prover\).
  \item \(\verifier\): \(S_\lambda\gets
        s_{u2}+\lambda s_{u4}\).
  \item \(\prover\): \(G_\lambda(X)\gets
        \widehat A(r_{u2},X)\MLE(u_2)(X)
        +\lambda\,\widehat A(r_{u4},X)\MLE(u_4)(X)
        +\lambda^2\chi_{r_{f1}}(X)\mathsf{TreeRes}(f_1,X)
        +\lambda^3\chi_{r_{g1}}(X)\mathsf{TreeRes}(g_1,X)
        +\lambda^4\chi_{r_{f2}}(X)\mathsf{TreeRes}(f_2,X)
        +\lambda^5\chi_{r_{g2}}(X)\mathsf{TreeRes}(g_2,X)\).
  \item \(\prover,\verifier\):
        \((a_{\mathrm{round}},r_{\mathrm{sc}},T_{\mathrm{sc}})
        \gets\SumCheck_n(\verifier:S_\lambda,\prover:G_\lambda)\).
  \item \(\verifier\): if \(a_{\mathrm{round}}=0\), return \((0,\emptyset)\).
  \item \(\prover\): let  \(w_{u2}\gets\MLE(u_2)(r_{\mathrm{sc}})\) and
        \(w_{u4}\gets\MLE(u_4)(r_{\mathrm{sc}}),\)\\
        \(\qquad w_{f1,p}\gets\MLE(f_1)(\mathsf{parent}(r_{\mathrm{sc}}))\),
        \(w_{f1,l}\gets\MLE(f_1)(\mathsf{left}(r_{\mathrm{sc}}))\),
        \(w_{f1,r}\gets\MLE(f_1)(\mathsf{right}(r_{\mathrm{sc}}))\),\\
        \(\qquad w_{g1,p}\gets\MLE(g_1)(\mathsf{parent}(r_{\mathrm{sc}}))\),
        \(w_{g1,l}\gets\MLE(g_1)(\mathsf{left}(r_{\mathrm{sc}}))\),
        \(w_{g1,r}\gets\MLE(g_1)(\mathsf{right}(r_{\mathrm{sc}}))\),\\
        \(\qquad w_{f2,p}\gets\MLE(f_2)(\mathsf{parent}(r_{\mathrm{sc}}))\),
        \(w_{f2,l}\gets\MLE(f_2)(\mathsf{left}(r_{\mathrm{sc}}))\),
        \(w_{f2,r}\gets\MLE(f_2)(\mathsf{right}(r_{\mathrm{sc}}))\),\\
        \(\qquad w_{g2,p}\gets\MLE(g_2)(\mathsf{parent}(r_{\mathrm{sc}}))\),
        \(w_{g2,l}\gets\MLE(g_2)(\mathsf{left}(r_{\mathrm{sc}}))\),
        \(w_{g2,r}\gets\MLE(g_2)(\mathsf{right}(r_{\mathrm{sc}}))\).
  \item \(\prover\to\verifier\): send \(w_{u2},w_{u4},\)\\
        \(\qquad w_{f1,p},w_{f1,l},w_{f1,r},w_{g1,p},w_{g1,l},w_{g1,r},\)\\
        \(\qquad w_{f2,p},w_{f2,l},w_{f2,r},w_{g2,p},w_{g2,l},w_{g2,r}\in\F\).  
  \item \(\verifier\): \(T_{\mathrm{endpoint}}\gets
        \widehat A(r_{u2},r_{\mathrm{sc}})
        w_{u2}
        +\lambda\,\widehat A(r_{u4},r_{\mathrm{sc}})
        w_{u4}
        +\lambda^2\chi_{r_{f1}}(r_{\mathrm{sc}})
        (w_{f1,p}-w_{f1,l}w_{f1,r})
        +\lambda^3\chi_{r_{g1}}(r_{\mathrm{sc}})
        (w_{g1,p}-w_{g1,l}w_{g1,r})
        +\lambda^4\chi_{r_{f2}}(r_{\mathrm{sc}})
        (w_{f2,p}-w_{f2,l}w_{f2,r})
        +\lambda^5\chi_{r_{g2}}(r_{\mathrm{sc}})
        (w_{g2,p}-w_{g2,l}w_{g2,r})\).
  \item \(\verifier\): if \(T_{\mathrm{sc}}\ne T_{\mathrm{endpoint}}\), return \((0,\emptyset)\).
  \item \(\verifier\): \(\EvalReqs_{\mathrm{sc}}\gets
        \{(\widehat O_{u_2},r_{\mathrm{sc}},w_{u2}),
        (\widehat O_{u_4},r_{\mathrm{sc}},w_{u4}),\)\\
        \(\qquad(\widehat O_{f_1},\mathsf{parent}(r_{\mathrm{sc}}),w_{f1,p})\),
        \((\widehat O_{f_1},\mathsf{left}(r_{\mathrm{sc}}),w_{f1,l})\),
        \((\widehat O_{f_1},\mathsf{right}(r_{\mathrm{sc}}),w_{f1,r})\),\\
        \(\qquad(\widehat O_{g_1},\mathsf{parent}(r_{\mathrm{sc}}),w_{g1,p})\),
        \((\widehat O_{g_1},\mathsf{left}(r_{\mathrm{sc}}),w_{g1,l})\),
        \((\widehat O_{g_1},\mathsf{right}(r_{\mathrm{sc}}),w_{g1,r})\),\\
        \(\qquad(\widehat O_{f_2},\mathsf{parent}(r_{\mathrm{sc}}),w_{f2,p})\),
        \((\widehat O_{f_2},\mathsf{left}(r_{\mathrm{sc}}),w_{f2,l})\),
        \((\widehat O_{f_2},\mathsf{right}(r_{\mathrm{sc}}),w_{f2,r})\),\\
        \(\qquad(\widehat O_{g_2},\mathsf{parent}(r_{\mathrm{sc}}),w_{g2,p})\),
        \((\widehat O_{g_2},\mathsf{left}(r_{\mathrm{sc}}),w_{g2,l})\),
        \((\widehat O_{g_2},\mathsf{right}(r_{\mathrm{sc}}),w_{g2,r})\}\)

  \item \(\verifier\): return \((1,\EvalReqs_{\mathrm{sc}})\).
\end{enumerate}
\end{protocolbox}

The terminal request set is explicit in Steps S14--S16: those are the values
that the later authentication layer must bind to the committed tensors.  The
first two requests are ordinary MLE endpoint claims for the accumulator inputs.
Each product-tree residual contributes three formal-extension endpoint claims:
parent, left child, and right child at the terminal sumcheck point
\(r_{\mathrm{sc}}\).  The third entry of each \(B\)-tuple is the handle used in
the endpoint request.  In the inlined Blaze protocol below, those entries are
hatted MLE handles such as \(\widehat O_{u_2}\) and \(\widehat O_{f_1}\).
The public factors \(\widehat A(r_{u2},r_{\mathrm{sc}})\),
\(\widehat A(r_{u4},r_{\mathrm{sc}})\), \(\chi_{r_{f1}}(r_{\mathrm{sc}})\),
\(\chi_{r_{g1}}(r_{\mathrm{sc}})\), \(\chi_{r_{f2}}(r_{\mathrm{sc}})\),
\(\chi_{r_{g2}}(r_{\mathrm{sc}})\), and the powers of \(\lambda\) are not oracle
requests.

\subsection{Blaze In The \(\mathcal F_{\mathrm{MLIOP}}\)-Hybrid}

We can now state the top-level Blaze flow as a hybrid protocol.  In this
subsection, \(\mathcal F_{\mathrm{MLIOP}}\) is the ideal multilinear-oracle
functionality from Figure~\ref{fig:func-mliop-oracle}.  This functionality is
the only oracle power Blaze needs at the native layer: the prover and verifier
name committed Boolean oracle handles, the functionality commits those handles
into MLE oracle handles, and the verifier later obtains MLE evaluations at
field points.

Figure~\ref{fig:native-blaze-fmliop-hybrid} is therefore not an incomplete
protocol and not merely a list of constraints.  It is Blaze in the
\(\mathcal F_{\mathrm{MLIOP}}\)-hybrid model.  The inlined RAA flow produces
evaluation requests of the form \(J(r)=v\), and the top-level protocol
immediately discharges those requests by calling
\(\mathcal F_{\mathrm{MLIOP}}.\mathsf{Eval}\).  Give this protocol real
implementations of the coordinate-oracle layer and of
\(\mathcal F_{\mathrm{MLIOP}}\), and the composition is a concrete Blaze
implementation.  Those implementations may use coordinate commitments,
batching, BaseFold, Merkle trees, or something else; the Blaze algebra below
only sees the hybrid functionality.

The important topology is the same one frozen in
Figure~\ref{fig:protocol-blaze-pcs}.  The encoded matrix is committed once as a
coordinate oracle; its rows are not separately committed to the MLIOP layer.
Only after the verifier samples \(\rho\) do the parties form the virtual folded
handle
\[
  R_\star(b)
  =
  \sum_{h\in\B^\ell}\rho[h]\,R_c(b)[h],
  \qquad b\in\B^n.
\]
This is the first outer encoded object passed to the inlined RAA flow.  The
protocol commits \(R_\star\) together with the other RAA-stage handles before
it asks any MLE-evaluation questions.
The folded tensor is
\[
  c_\star[b]=\sum_{h\in\B^\ell}\rho[h]c[h,b],
  \qquad b\in\B^n.
\]

\begin{protocolbox}{Protocol: Blaze In The \(\mathcal F_{\mathrm{MLIOP}}\)-Hybrid}{fig:native-blaze-fmliop-hybrid}
\scriptsize\raggedright
\noindent Public parameters: \(\ell,k,n\), the RAA map
\(\RAA:\F^{\B^k}\to\F^{\B^n}\), \(\pi_1,\pi_2:\B^n\to\B^n\),
\(\widehat\pi_1,\widehat\pi_2,\Idx:\F^n\to\F\),
\(\widehat A:\F^n\times\F^n\to\F\), \(F_{\mathrm{rep}}\),
\(\mathsf{msg}:\F^n\to\F^k\), \(\mathsf{RowFold}\),
\(\AccReduce\), \(\PermuteCommit\), \(\PermuteReduce\), \(\RAASixSumcheck\),
\(\mathcal F_{\mathrm{MsgOracle}}\), and \(\mathcal F_{\mathrm{MLIOP}}\).

\smallskip
\noindent\(\underline{\mathsf{Blaze}^{\mathcal F_{\mathrm{MLIOP}}}.\mathsf{Commit}(\verifier:\bot,\ \prover:m\in\F^{\B^\ell\times\B^k})\to C_m\in\Handle}\)

\smallskip
\begin{enumerate}[leftmargin=0.24in,label=\textbf{C\arabic*.},itemsep=0.06em,topsep=0.06em,parsep=0pt]
  \item \(\prover\): for each \(h\in\B^\ell\), compute
        \(c[h,\ast]\gets\RAA(m[h,\ast])\in\F^{\B^n}\).
  \item \(\prover\): define \(C\in(\F^{\B^\ell})^{\B^n}\) by
        \(C[b]=c[\ast,b]\) for \(b\in\B^n\).
  \item \(\prover,\verifier\): \(R_c\gets
        \mathcal F_{\mathrm{MsgOracle}}.\mathsf{Submit}
        (\prover:C,\verifier:\bot)\).
  \item \(\mathsf{Blaze}^{\mathcal F_{\mathrm{MLIOP}}}\): return
        \(C_m=R_c\).
\end{enumerate}

\smallskip
\noindent\(\underline{\mathsf{Blaze}^{\mathcal F_{\mathrm{MLIOP}}}.\mathsf{Open}(\verifier:(C_m,z_{\mathrm{slice}}\in\F^\ell,z_{\mathrm{in}}\in\F^k,v\in\F),\ \prover:(m,C_m,z_{\mathrm{slice}},z_{\mathrm{in}},v))\to a\in\{0,1\}}\)

\smallskip
\begin{enumerate}[leftmargin=0.24in,label=\textbf{O\arabic*.},itemsep=0.06em,topsep=0.06em,parsep=0pt]
  \item \(\prover\): for each \(h\in\B^\ell\), compute
        \(u[h]\gets\MLE(m[h,\ast])(z_{\mathrm{in}})\).
  \item \(\prover\to\verifier\): \(u\in\F^{\B^\ell}\).
  \item \(\verifier\): if
        \(v\ne\sum_{h\in\B^\ell}\chi_h(z_{\mathrm{slice}})u[h]\), return
        \(0\).
  \item \(\verifier\): sample \(\rho\leftarrow\F^{\B^\ell}\), then send
        \(\rho\) to \(\prover\).
  \item \(\prover\): compute
        \(m_\star[a]\gets\sum_{h\in\B^\ell}\rho[h]m[h,a]\) for
        \(a\in\B^k\).
  \item \(\verifier\): let
        \(v_\star\gets\sum_{h\in\B^\ell}\rho[h]u[h]\).
  \item \(\prover,\verifier\): \(R_\star\gets
        \mathsf{RowFold.Commit}(\verifier:(C_m,\rho),\prover:(C_m,\rho))\).
  \item \(\prover\): let \(u_0\gets m_\star\),
        \(u_1\gets F_{\mathrm{rep}}(u_0)\),
        \(u_2\gets M_{\pi_1}u_1\),
        \(u_3\gets\Acc(u_2)\),
        \(u_4\gets M_{\pi_2}u_3\),
        \(c_\star^{\mathrm{wit}}\gets\Acc(u_4)\), and
        \(u_5\gets c_\star^{\mathrm{wit}}\).
  \item \(\prover,\verifier\): let \(R_{u_5}\gets R_\star\).
  \item \(\prover,\verifier\): for each \(i\in\{0,2,3,4\}\), let
        \(R_{u_i}\gets
        \mathcal F_{\mathrm{MsgOracle}}.\mathsf{Submit}(\prover:u_i,\verifier:\bot)\).
  \item \(\prover,\verifier\): for each \(i\in\{0,2,3,4,5\}\), let
        \(\widehat O_{u_i}\gets
        \mathcal F_{\mathrm{MLIOP}}.\mathsf{Commit}(\prover,\verifier:R_{u_i})\).
  \item \(\prover,\verifier\): let
        \(\widehat O_{u_1}\gets\RepView(\widehat O_{u_0})\).
  \item \(\verifier\): sample \(\alpha,\beta\leftarrow\F\), then send
        \((\alpha,\beta)\) to \(\prover\).
  \item \(\prover,\verifier\):
        \((\widehat O_{f_1},\widehat O_{g_1};f_1,g_1)\gets
        \PermuteCommit(\verifier:(\widehat O_{u_1},
        \widehat O_{u_2},\alpha,\beta),\prover:(u_1,u_2))\).
  \item \(\prover,\verifier\):
        \((\widehat O_{f_2},\widehat O_{g_2};f_2,g_2)\gets
        \PermuteCommit(\verifier:(\widehat O_{u_3},
        \widehat O_{u_4},\alpha,\beta),\prover:(u_3,u_4))\).
  \item \(\verifier\): sample \(r_{\mathrm{chk}}\leftarrow\F^n\), then send
        \(r_{\mathrm{chk}}\) to \(\prover\).
  \item \(\prover,\verifier\):
        \((\EvalReqs_{u3},B_{u2})\gets
        \AccReduce(\verifier:(\widehat O_{u_2},
        \widehat O_{u_3},r_{\mathrm{chk}}),\prover:(u_2,u_3))\).
  \item \(\prover,\verifier\):
        \((\EvalReqs_{y},B_{u4})\gets
        \AccReduce(\verifier:(\widehat O_{u_4},
        \widehat O_{u_5},r_{\mathrm{chk}}),\prover:(u_4,u_5))\).
  \item \(\prover,\verifier\):
        \((\EvalReqs_{f1,g1},B_{f1},B_{g1})\gets
        \PermuteReduce(\verifier:(\widehat O_{u_1},
        \widehat O_{u_2},\widehat O_{f_1},\widehat O_{g_1},
        \widehat\pi_1,\alpha,\beta,r_{\mathrm{chk}}),
        \prover:(u_1,u_2,f_1,g_1))\).
  \item \(\prover,\verifier\):
        \((\EvalReqs_{f2,g2},B_{f2},B_{g2})\gets
        \PermuteReduce(\verifier:(\widehat O_{u_3},
        \widehat O_{u_4},\widehat O_{f_2},\widehat O_{g_2},
        \widehat\pi_2,\alpha,\beta,r_{\mathrm{chk}}),
        \prover:(u_3,u_4,f_2,g_2))\).
  \item \(\prover,\verifier\): \((a_{\mathrm{sc}},\EvalReqs_{\mathrm{sc}})\gets
        \RAASixSumcheck(\verifier:(B_{u2},B_{u4},
        B_{f1},B_{g1},B_{f2},B_{g2}),
        \prover:(u_2,u_4,f_1,g_1,f_2,g_2))\).
        \item \(\verifier\): \(\EvalReqs_{\RAA}\gets
        \{(\widehat O_{u_0},z_{\mathrm{in}},v_\star)\}
        \cup\EvalReqs_{u3}\cup\EvalReqs_y
        \cup\EvalReqs_{f1,g1}\cup\EvalReqs_{f2,g2}
        \cup\EvalReqs_{\mathrm{sc}}\).
        \item \(\verifier\): for each \((J,r,w)\in\EvalReqs_{\RAA}\): \(\verifier\): if \( \mathcal F_{\mathrm{MLIOP}}.\mathsf{Eval}
          (\verifier:(J,r))\ne w\), return \(0\).
        
        \item \(\verifier\): if \(a_{\mathrm{sc}}=0\), return \(0\).
  \item \(\verifier\): return \(1\).
\end{enumerate}
\end{protocolbox}

Figure~\ref{fig:native-blaze-fmliop-hybrid} is the promised Blaze protocol in
the hybrid model.  It uses coordinate handles, but not Merkle trees or BaseFold
query paths.  Those are possible ways to realize the two ideal layers:
\(\mathcal F_{\mathrm{MsgOracle}}\) for coordinate access and
\(\mathcal F_{\mathrm{MLIOP}}\) for MLE evaluation access.  The next layer is
therefore not a second RAA protocol.  It is an implementation of the ideal
calls used by this figure.

The Blaze paper uses a different presentation style.  It first constructs an
MLIOP for \(\RMLE[\PRAA]\), then describes an MLIOP-to-IOPP transformation using
encoded oracle strings, batching, and an IOPP for \(\RMLE[C]\)~[Blaze24,
Section~4.1].  In the language of this explainer, that transformation is a
realization of \(\mathcal F_{\mathrm{MLIOP}}\).  The following dictionary is the
translation:
\[
\begin{array}{c|c}
\text{this explainer} & \text{Blaze paper} \\
\hline
\mathcal F_{\mathrm{MLIOP}}\text{-hybrid} &
\text{MLIOP model} \\
\widehat O &
\text{multilinear oracle } \widehat H \\
\mathcal F_{\mathrm{MLIOP}}.\mathsf{Commit}(R) &
\text{ideal wrapper for authenticated access to }\widehat H \\
\left(\widehat O,r,v\right)\in\EvalReqs &
\text{verifier evaluation query } \widehat H(r)=v \\
\text{realize }\mathcal F_{\mathrm{MLIOP}} &
\text{encoded-oracle/RMLE realization via the MLIOP-to-IOPP transformation} \\
\MLEAuth &
\Pi_{\mathrm{batch}}\text{ plus an IOPP for }\RMLE[C] \\
\text{virtual handle} &
\text{virtual multilinear polynomial/oracle}
\end{array}
\]

\subsection{Boundary To The Compiled Protocol}

The native section stops at
Figure~\ref{fig:native-blaze-fmliop-hybrid}.  It has checked the native
sumcheck transcript and all requested MLE evaluations inside the
\(\mathcal F_{\mathrm{MLIOP}}\)-hybrid model.  A concrete protocol replaces the
hybrid functionality by a real implementation.  The implementation described
later realizes the needed Boolean-handle commits and MLE-evaluation calls using
\(\MLEAuth\), \(\BackendRMLE\), BaseFold-style proximity, and coordinate
openings.

This is why \(\EvalReqs\) remain the right internal boundary object for the RAA
algebra.  A native MLIOP request says only, ``the verifier needs
\(J(r)=v\).''  In the hybrid protocol, the final evaluation loop checks that
request by calling \(\mathcal F_{\mathrm{MLIOP}}\).  In a concrete implementation, the
realization of \(\mathcal F_{\mathrm{MLIOP}}\) adds coordinate-opening handles,
batches such requests, and uses an RMLE backend to authenticate the batched MLE
claim.  The RAA algebra does not need to know whether that handle is
ordinary, composite, or virtual.

This separation is the main invariant to preserve:
\[
  \text{RAA algebra creates evaluation requests;}
  \qquad
  \text{the backend authenticates them.}
\]
